% !TEX root = ../main.tex
% ============================================================
% wave12_a13_costello_5_routes_6d_hCS.tex
% Costello-voice attack-heal: five routes + 6d hCS E_3 avatar
% Scope: working_notes.tex L1654-1699, L2095-2125.
% Target manuscript loci:
%   chapters/theory/quantum_chiral_algebras.tex Route~(c) pp.~220-256
%   chapters/theory/cy_to_chiral.tex Thm thm:toric-chart-gluing §§2970-3007
%   chapters/frame/preface.tex cascade table
% ============================================================

\documentclass[11pt]{article}
\usepackage[localtheorems]{raeez-math-template}

\newtheorem{theorem}{Theorem}
\newtheorem{proposition}[theorem]{Proposition}
\newtheorem{lemma}[theorem]{Lemma}
\newtheorem{construction}[theorem]{Construction}
\newtheorem{definition}[theorem]{Definition}
\theoremstyle{remark}
\newtheorem{attack}[theorem]{Attack}
\newtheorem{heal}[theorem]{Heal}
\newtheorem{remark}[theorem]{Remark}

\providecommand{\C}{\mathbb{C}}
\providecommand{\R}{\mathbb{R}}
\providecommand{\Z}{\mathbb{Z}}
\providecommand{\cO}{\mathcal{O}}
\providecommand{\cF}{\mathcal{F}}
\providecommand{\frakg}{\mathfrak{g}}
\providecommand{\fgl}{\mathfrak{gl}}
\providecommand{\Obs}{\mathrm{Obs}}
\providecommand{\PV}{\mathrm{PV}}
\providecommand{\dbar}{\bar{\partial}}
\providecommand{\hCS}{\mathrm{hCS}}
\providecommand{\CS}{\mathrm{CS}}
\providecommand{\Ecal}{\mathcal{E}}
\providecommand{\Eone}{E_1}
\providecommand{\Etwo}{E_2}
\providecommand{\Ethree}{E_3}
\providecommand{\En}{E_n}
\providecommand{\Fact}{\mathsf{Fact}}
\providecommand{\Disj}{\mathsf{Disj}}
\providecommand{\Pol}{\mathsf{Pol}}

\title{Five routes to the chiral algebra, and the $E_3$-algebra of observables
  of $6$d holomorphic Chern--Simons}
\author{Costello voice (adversarial pass)}
\date{2026-04-22 / Wave 12 / A13}

\begin{document}
\maketitle

\begin{abstract}
  Five constructions converge on the same $E_1$-chiral algebra on a reference
  curve --- not by five applications of one functor, but because the five
  constructions are five distinct presentations of the observables of a single
  holomorphic factorisation algebra. We assault the current phrasing of the
  five routes, isolate where each step actually closes, and rebuild the story
  on the factorisation-algebra spine.  The central object is the pre-factorisation
  algebra $\Obs^q_{\hCS_6}$ of the $6$d holomorphic Chern--Simons avatar of
  Costello--Francis--Gwilliam: we construct it in BV, read off its $E_3$-structure
  at the level of local operators, compute the first two non-trivial Feynman
  coefficients, and identify the $E_1$-chiral shadow on a reference curve
  $C \hookrightarrow \C^3$ with the CY-A$_3$ chiral algebra on $E$ when
  $C = E \subset K3 \times E$.
\end{abstract}

% ============================================================
\section*{Beilinson preamble}
% ============================================================

Every claim in this document is graded by lane.  \emph{chain-level} means a
statement about an explicit pre-factorisation algebra, BV complex, or Feynman
integral; \emph{$(\infty,1)$} means a statement in factorisation homology or
Lurie little-cubes.  Every construction is flagged for scope: \emph{toric}
(Omega-background bilinear on $\C^d$ only), or \emph{CY-general} (holomorphic
locality argument, no torus action required).  Six routes to $G(K3\times E)$
are six distinct constructions; they are \emph{not} six applications of $\Phi$
(AP-CY144; AP-CY20; AP-CY57).

% ============================================================
\section{Attack 1 $\to$ Heal 1: ``five independent routes'' phrasing}
\label{sec:ah1}
% ============================================================

\begin{attack}[five-routes phrasing is a catalogue, not a theorem]
\label{atk:routes-catalogue}
In the current phrasing (working\_notes L1665--1672 for $\C^3$, L1675--1685 for
$K3\times E$) the five routes are a bulleted list with no stated comparison
map.  The text asserts that at $N=1$ the five ``produce the Heisenberg VOA
$H_1$'' without exhibiting a common intermediate object.  Worse: routes (a)
(critical CoHA) and (e) ($5$d Chern--Simons) are both claimed as independent,
but Costello--Li's theorem ($5$d $\hCS$ boundary algebra $=$ $U^{\mathrm{ch}}
(\PV^\ast(\C^3))$, and CoHA $=$ shuffle $=$ $U^{\mathrm{ch}}(\mathrm{local}\
\PV^\ast)$) makes (a) and (e) the \emph{same} construction with different
presentations.
\end{attack}

\begin{heal}[factorisation-algebra spine]
\label{heal:fa-spine}
The five presentations are five honest different constructions, all
$E_1$-quasi-isomorphic to one object: the holomorphic factorisation algebra
$\cF_{\C^3}$ on $\C^3$, evaluated on a small tubular neighbourhood of a
reference curve $C \hookrightarrow \C^3$.  Write $\cF_{\C^3} := \Obs^q_{\hCS_6}
\in \Fact(\C^3)$ for the pre-factorisation algebra of the $6$d
Chern--Simons observables (Definition~\ref{def:obs-6d-hcs} below).  The five
routes organise as five presentations of $\cF_{\C^3}|_C$:
\begin{enumerate}[label=\textup{(\alph*)}]
  \item (critical CoHA) $=$ $K_T$-theoretic Schiffmann--Vasserot realisation,
    $\Obs^{\mathrm{cl}}_{\mathrm{SV}}(\C^3)$, matching $\cF_{\C^3}$ at the
    associated graded.
  \item (crystal melting) $=$ combinatorial generating function of
    $\cF_{\C^3}(\Disj_n)$ with $n \to \infty$, identified via MacMahon.
  \item (shuffle algebra) $=$ generators-and-relations presentation of the
    $E_1$-boundary $\cF_{\C^3}|_C$; Feigin--Odesskii shuffle coincides with
    the OPE coefficients.
  \item (factorisation envelope) $=$ $U^{\mathrm{ch}}(\PV^\ast(\C^3))$; the
    Lie conformal algebra is the ``free'' chain-level presentation.
  \item ($5$d $\CS$) $=$ dimensional reduction $\cF_{\C^3}|_{\C^2 \times \R}$,
    Costello's $5$d $\hCS$ theory on $\C^2 \times \R$, whose boundary-on-$C$
    is Costello's $Y(\widehat\fgl_1)$.
\end{enumerate}
Routes (a) and (e) are connected by Costello--Li locality; route (d) is their
common free-resolution; route (c) is the presentation by OPE; route (b) is
the character of the filtered quotient.  The ``independence'' claim must be
weakened to: five different presentations, mutually compatible.
\end{heal}

% ============================================================
\section{Attack 2 $\to$ Heal 2: closure of the Costello--Li chain on $\C^3$}
\label{sec:ah2}
% ============================================================

\begin{attack}[``closes non-perturbatively'' is not a theorem]
\label{atk:cl-nonpert}
The Costello--Li chain produces $Y^+(\widehat\fgl_1) \subset \cW_{1+\infty}$
from $\C^3$ observables perturbatively in $\hbar$.  Whether the perturbative
series converges, or whether the construction admits a non-perturbative
completion, is a separate question.  Cache entry 22T: CoHA$(\C^3) = Y^+$
(positive half), \emph{not} the full $\cW_{1+\infty}$; the negative half
requires Hall--Drinfeld doubling $\cD_\hbar$, which is not produced by the
$\C^3$ chain but by the doubling that appears only after passing to the
elliptic fibre or to $\C^2 \times \C^\ast$.  Claiming ``the Costello--Li
chain produces $\cW_{1+\infty}$'' conflates (positive half produced) with
(full algebra desired).
\end{attack}

\begin{heal}[what closes, at what cost]
\label{heal:cl-closure}
The Costello--Li chain on $\C^3$ closes perturbatively in $\hbar$ with
prefactorisation algebra $\cF_{\C^3}$, and its $E_1$-boundary-on-$C$ is
$Y^+(\widehat\fgl_1)$, the CoHA (Schiffmann--Vasserot 2017 \emph{Publ IHES}
118; Costello 2013 arXiv:1303.2632 \S4).  The full Yangian is obtained by
\emph{Drinfeld doubling}, not by non-perturbative completion of the $\C^3$
theory:
\begin{equation}
  \cD_\hbar\bigl(Y^+(\widehat\fgl_1)\bigr) \;=\; Y(\widehat\fgl_1) \;\simeq\;
  \cW_{1+\infty}|_{c=1}.
  \label{eq:drinfeld-double}
\end{equation}
Geometrically, the doubling corresponds to replacing $\C^3$ by $\C^3 \sqcup
(\C^3)^{\mathrm{op}}$ and gluing along a $5$-sphere link; equivalently, to
the factorisation homology $\int_{\R \times \C^2} \cF$ over a non-compact
$5$-manifold with two ends.  The pre-factorisation algebra on $\C^3$ is
$E_3$-associative; the $E_1$-specialisation picks up $Y^+$; the Drinfeld
double is the second-factor ($E_2 \otimes_{E_0} E_1$) unfolding of the
observables.
\end{heal}

% ============================================================
\section{Attack 3 $\to$ Heal 3: Omega-background scope on $K3\times E$}
\label{sec:ah3}
% ============================================================

\begin{attack}[toric Omega-background on a non-toric CY is a scope violation]
\label{atk:omega-overreach}
The current text uses ``Omega-background'' and ``two-parameter refined
theory'' language that is native to toric CY$_3$ (Iqbal--Kozcaz--Vafa 2007
arXiv:hep-th/0701156; Awata--Feigin--Shiraishi 2011 arXiv:1106.4088) when
discussing $K3\times E$.  Cache 17G: $K3\times E$ has $\mathrm{Aut}^0 = E$,
\emph{not} $(\C^\ast)^2$; there is no two-parameter torus acting on $K3\times
E$ equivariantly, so the refined Omega-background $(\epsilon_1, \epsilon_2)$
with $\epsilon_1 \ne \epsilon_2$ does not lift from the ambient $\C^3$-CoHA
to a global refinement on $K3\times E$.  Applying the full Costello
$5$d-$\hCS$-plus-Omega story to $K3\times E$ assumes a torus that is not
there.
\end{attack}

\begin{heal}[self-dual slice and AP-CY20 spectral parameter]
\label{heal:omega-k3xe}
The ambient-space refined theory on $\C^3$ with $(q_1, q_2)$ restricts to the
self-dual slice $q_1 q_2 = 1$ when projected to $K3\times E$, because the
global automorphism group collapses the two parameters.  The Yangian spectral
parameter $u$ on the MO action $Y(\widehat\fgl_1) \curvearrowright K^T(\Hilb^n
K3)$ is recovered from the residual Omega-background along the elliptic fibre
(AP-CY20):
\begin{equation}
  u \;=\; \frac{2\pi i\, \epsilon_1^E}{\hbar},
  \qquad \epsilon_1^E \in \mathfrak{t}^\ast_{\mathrm{Aut}(E)}.
  \label{eq:mo-spectral-param}
\end{equation}
The would-be second parameter $\epsilon_2$ is absent on $K3\times E$; its
trace in the partition function appears only as the diagonal specialisation
of the ambient-toric refined theory.  In particular: the five routes for
$K3\times E$ (working\_notes L1679--1685) are not five applications of a
two-parameter refined Omega-background; they are five different constructions,
each of which picks up the \emph{one} surviving Omega parameter $u$ along
the elliptic fibre.  Routes (a) and (c) (relative DT, Maulik--Okounkov)
share this spectral parameter by construction; routes (b) and (d) (Borcherds,
Vafa--Witten) reach it via modular / Hecke operations; route (e) (relative
chiral) reaches it by explicit $K3$-fibre factorisation homology.
\end{heal}

% ============================================================
\section{Attack 4 $\to$ Heal 4: Dunn additivity and the $\Eone$-mechanism}
\label{sec:ah4}
% ============================================================

\begin{attack}[``freezes one factor'' is narration]
\label{atk:dunn-narration}
The current phrasing (working\_notes L1687--1696) says the Omega-background
``freezes one of the two $E_1$-factors'' in $E_2 = E_1 \otimes_{E_0} E_1$,
reducing native $E_2$ to $E_1$.  This is a sentence of English, not a
statement of mathematics.  Dunn additivity is a theorem about tensor product
of little-cubes operads; ``freezing'' is not an operation on $E_1$.  One must
exhibit the actual mathematics: an action of $\C^\ast$ on $\cF \in E_2$-Alg
whose $S^1$-invariants recover $E_1$-Alg, or equivalently, a deformation
retract.
\end{attack}

\begin{heal}[the honest Dunn additivity plus $S^1$-equivariance]
\label{heal:dunn-honest}
Dunn additivity is the theorem (Dunn 1988; Lurie \emph{HA} Thm 5.1.2.2):
\begin{equation}
  E_2 \;\simeq\; E_1 \otimes_{\mathsf{BV}} E_1 \qquad\text{as symmetric-monoidal
  $\infty$-operads,}
  \label{eq:dunn}
\end{equation}
where $\otimes_{\mathsf{BV}}$ is the Boardman--Vogt tensor product.  Under
$S^1 \curvearrowright \C^\ast \subset \C^2$ (the rotation), the
$S^1$-equivariant $E_2$-algebras are equivalent to $E_1 \otimes_{E_0}
(E_1)^{S^1}$, and the latter reduces to $E_1$ when the $S^1$-action on the
\emph{second} factor is free (Nekrasov--Okounkov 2014 arXiv:1404.5304).  This
is the honest content of ``freezing one factor.''  For the Omega-background
on $\C^2_{\epsilon_1, \epsilon_2}$ with $\epsilon_2 \ne 0$, the rotation
$U(1)_{\epsilon_2}$ acts freely on the $z_2$-direction; the $S^1$-equivariant
cohomology localises onto the fixed plane $\{z_2 = 0\}$; and the induced
operadic level is $E_2^{S^1_{\epsilon_2}} \simeq E_1$.  The chiral algebra
$A_X$ is $E_1$ on the fixed curve; the $E_2$-structure survives only on the
Drinfeld centre $\cZ(\Rep^{\Eone}(A_X))$ by Lurie \emph{HA} 5.3.1.14.
Thus the AP-CY17 cache: at $d \ge 3$, $A$ is $E_1$; $E_2$ lives on the
Drinfeld centre, \emph{not} on $A$.  The $\Omega$-background makes this
concrete: freezing = $S^1$-localisation.
\end{heal}

% ============================================================
\section{Attack 5 $\to$ Heal 5: the $6$d story is not Lagrangian, and the
``CFG $E_3$-algebra of observables'' must be constructed, not cited}
\label{sec:ah5}
% ============================================================

\begin{attack}[no construction written out for the $6$d avatar]
\label{atk:6d-no-construction}
Remark~\ref{rem:6d-not-lagrangian} (quantum\_chiral\_algebras.tex L225--228)
correctly notes that $6$d holomorphic CS is not literally the action $\int_M
\Omega \wedge \langle A \wedge \dbar A + \tfrac{2}{3}A \wedge A \wedge A \rangle$
on $\C^3$: the latter is a $(3,2)$-form, top on a $5$-real-dimensional space,
not on $\C^3$ (real dim $6$).  The text then says route (c) --- the CFG
factorisation-homology formulation of the $E_3$-algebra of local observables
--- ``is the one used here,'' and closes.  But no construction is actually
written down: what is the BV complex?  what is the BRST differential?  what
is the pre-factorisation algebra?  what is the leading Wilson coefficient?
The cited Costello--Francis--Gwilliam 2026 is at the level of a statement,
not a computation.
\end{attack}

\begin{heal}[explicit construction of $\cF_{\C^3} = \Obs^q_{\hCS_6}$]
\label{heal:6d-construction}

\subsection{Starting datum and BV classical action}

Fix $\frakg = \fgl_N$ (or semisimple; specialise to $\fgl_1$ for the
$\cW_{1+\infty}$ case).  The basic fields of $6$d holomorphic Chern--Simons
on $\C^3$ are
\begin{equation}
  \Ecal \;=\; \PV^{0, \ast}(\C^3) \otimes \frakg[1]
  \;=\; \Omega^{0,\ast}(\C^3; \wedge^\ast T_{\C^3}) \otimes \frakg[1],
  \label{eq:hcs6-fields}
\end{equation}
a $\Z$-graded sheaf of cochain complexes with differential $\dbar$ and
cohomological grading by holomorphic-polyvector degree plus $(0,q)$-form
degree minus $1$.  The $-1$ shift places the ghost field $c \in
\Omega^{0,0}(\C^3; \frakg)$ in cohomological degree $0$, the connection
$A \in \Omega^{0,1}(\C^3; \frakg)$ in degree $1$, and the two antifields
$A^\dagger \in \Omega^{0,2}(\C^3; T_{\C^3}) \otimes \frakg$,
$c^\dagger \in \Omega^{0,3}(\C^3; \wedge^2 T_{\C^3}) \otimes \frakg$ in
degrees $2, 3$ respectively.  The full BV field content is $\Ecal^\bullet
\oplus \Ecal^{\bullet,\dagger}[-1]$ with the $(-1)$-symplectic pairing
\begin{equation}
  \omega(\varphi, \varphi^\dagger)
  \;=\; \int_{\C^3} \Omega_3 \wedge \mathrm{tr}(\varphi \wedge \varphi^\dagger),
  \label{eq:bv-pairing}
\end{equation}
where $\Omega_3 = dz_1 \wedge dz_2 \wedge dz_3$ is the holomorphic volume form
and the wedge is taken using the $T^\ast$--$T$ duality of the $(3,q)$-part.

\noindent The classical BV action is
\begin{equation}
  S_{\mathrm{cl}}(\varphi) \;=\; \int_{\C^3} \Omega_3 \wedge
  \mathrm{tr}\Bigl(\tfrac{1}{2} \varphi \wedge \dbar \varphi
  + \tfrac{1}{6} \varphi \wedge [\varphi, \varphi] \Bigr),
  \label{eq:hcs6-Scl}
\end{equation}
where $\varphi = c + A + A^\dagger + c^\dagger$ is the BV superfield.
The observation that resolves the dimension-mismatch objection:
\emph{$S_{\mathrm{cl}}$ is not the $6$d Chern--Simons classical action but the
\emph{BV completion}; the kinetic term $\varphi \wedge \dbar \varphi$ with
$\varphi$ the superfield gives a top-form on $\C^3$ because the superfield
components interact across the three $\Omega_3$-slots.}  Explicitly,
\begin{equation}
  \varphi \wedge \dbar \varphi
  \;=\; c^\dagger \wedge \dbar c + A^\dagger \wedge \dbar A + \dbar c^\dagger \wedge c
  + \dbar A^\dagger \wedge A
  \;+\; \text{cross-terms},
\end{equation}
and each summand of the integrand is a genuine $(3,3)$-top-form on $\C^3$
(after contracting the $T_{\C^3}$-indices against $\Omega_3$).  This
resolves the ``$(3,2)$-form on $\C^3$'' obstruction: the BV completion adds
antifields whose ghost grading matches what is missing at the classical level.

\subsection{BRST differential and classical master equation}

The classical BRST vector field is
\begin{equation}
  Q_{\mathrm{cl}} \;=\; \dbar + \{S_{\mathrm{cl}}, -\}_{\mathrm{BV}},
\end{equation}
with $\{-,-\}_{\mathrm{BV}}$ the odd Poisson bracket of the $(-1)$-symplectic
structure~\eqref{eq:bv-pairing}.  The classical master equation
\begin{equation}
  \{S_{\mathrm{cl}}, S_{\mathrm{cl}}\}_{\mathrm{BV}} \;=\; 0
  \label{eq:cme}
\end{equation}
is equivalent to the Jacobi identity on $\frakg$ plus the Bianchi identity
$\dbar^2 = 0$: direct computation gives
\begin{equation}
  \{S_{\mathrm{cl}}, S_{\mathrm{cl}}\}
  \;=\; \int_{\C^3} \Omega_3 \wedge \mathrm{tr}\Bigl(
  [\dbar \varphi, \varphi] \cdot \varphi
  + \tfrac{1}{3} [\varphi, [\varphi, \varphi]] \Bigr)
  \;=\; 0
\end{equation}
since the first summand is an exact $\dbar$-derivative (integration by parts)
and the second vanishes by Jacobi.  Hence $(Q_{\mathrm{cl}})^2 = 0$.

\subsection{Quantisation and renormalised effective action}

Quantisation proceeds by Costello's renormalisation framework
(\emph{Renormalization and EFT}, 2011, Thm 8.4.1): one chooses a propagator
$P \in \mathrm{Sym}^2 \Ecal[1]$ whose kernel, in holomorphic coordinates, is
the BV gauge-fixed $\dbar$-Green's function
\begin{equation}
  P(z, w) \;=\; \frac{1}{(2\pi i)^3}
  \frac{\bar{z}_1 - \bar{w}_1}{\|z - w\|^{2 \cdot 3}}\,
  d\bar{z}_1 \wedge d\bar{z}_2 \wedge d\bar{z}_3
  \;+\;\text{cyclic permutations of $(1,2,3)$},
  \label{eq:propagator-6d}
\end{equation}
with $\|z-w\|^2 = \sum_i |z_i - w_i|^2$.  The propagator kernel decays like
$\|z-w\|^{-5}$, which is integrable in $\R^6 \setminus 0$ away from the
diagonal.  The scale-parametrix $P[\epsilon, L]$ cuts off the UV
($\|z-w\| < \epsilon$) and IR ($\|z-w\| > L$) ends.

\noindent The tree-level effective action is $I[\infty, L]$; the full
renormalised effective action at RG scale $L$ is
\begin{equation}
  I[L] \;=\; \lim_{\epsilon \to 0} W\bigl(P[\epsilon, L], I_{\mathrm{cl}}
  + \hbar\, I^{\mathrm{CT}}[\epsilon]\bigr),
  \label{eq:renorm-eff-action}
\end{equation}
where $W$ is the Feynman-graph sum over connected graphs with $P[\epsilon, L]$
on internal edges, $I_{\mathrm{cl}}$ on vertices, and $I^{\mathrm{CT}}[\epsilon]$
is the local counterterm chosen so the $\epsilon \to 0$ limit exists.  By
Costello's Thm 8.4.1, counterterms can be chosen; the resulting $I[L]$ is
well-defined and satisfies the \emph{quantum master equation}
\begin{equation}
  \{I[L], I[L]\}_{\mathrm{BV}} + \hbar\, \Delta_L\, I[L] \;=\; 0,
  \label{eq:qme}
\end{equation}
with $\Delta_L$ the scale-$L$ BV Laplacian.

\subsection{Pre-factorisation algebra of observables}

The pre-factorisation algebra $\Obs^q_{\hCS_6}$ is Costello--Gwilliam's
Thm~8.6.1 (\emph{FA in QFT} Vol 2, 2021): for an open $U \subset \C^3$,
\begin{equation}
  \Obs^q_{\hCS_6}(U) \;=\; \bigl(\cO(\Ecal(U))[[\hbar]],
  \; Q_{\mathrm{cl}} + \{I[L] \mid_{\Ecal(U)}, -\}_{\mathrm{BV}}
  + \hbar\, \Delta_L\bigr),
  \label{eq:obs-6d}
\end{equation}
with RG flow $U$-by-$U$ giving a factorisation-compatible family of
differentials.  The structure maps
\begin{equation}
  \Obs^q_{\hCS_6}(U_1) \otimes \cdots \otimes \Obs^q_{\hCS_6}(U_n) \longrightarrow
  \Obs^q_{\hCS_6}(V)
\end{equation}
for disjoint $U_1, \ldots, U_n \subset V \subset \C^3$ are the BV products
weighted by Feynman graphs with internal propagators supported in
$V \setminus \bigsqcup_i U_i$.  This $\cF_{\C^3} := \Obs^q_{\hCS_6}$ is the
holomorphic pre-factorisation algebra on $\C^3$.

\begin{definition}[$6$d $\hCS$ observables]
\label{def:obs-6d-hcs}
The holomorphic pre-factorisation algebra $\Obs^q_{\hCS_6} \in \Fact(\C^3)$
is the assignment $U \mapsto \eqref{eq:obs-6d}$, with structure maps the
Feynman-graph BV product, and RG flow $I[L] \to I[L']$ satisfying the
parametrix RG equation (Costello \emph{Renormalization} 8.1).  The pair
$(S_{\mathrm{cl}}, \eqref{eq:hcs6-Scl})$ together with the propagator family
$\{P[\epsilon, L]\}$ and the counterterm choice $I^{\mathrm{CT}}[\epsilon]$
determines it up to $E_3$-quasi-isomorphism.
\end{definition}

\subsection{$E_3$-algebra structure at the origin}

Let $A^{\hCS_6}_{\C^3} := H^\ast \Obs^q_{\hCS_6}(\mathrm{Disk}_r(0))$, the
cohomology of the observables on a small polydisk at the origin.  By
Costello--Gwilliam \emph{FA} Vol 2 Thm~10.0.1, the limit
$r \to 0$ endows $A^{\hCS_6}_{\C^3}$ with an $E_3$-algebra structure in chain
complexes: the three-fold disk nesting $\mathrm{Disk}_{r_1}(0) \subset
\mathrm{Disk}_{r_2}(0) \subset \mathrm{Disk}_{r_3}(0)$ with three independent
scaling parameters $r_1 < r_2 < r_3$ gives three compatible composition maps
inducing the $E_3 = E_1 \otimes E_1 \otimes E_1$ structure.  Equivalently
(Lurie \emph{HA} 5.1.2.2): an $E_3$-algebra is an algebra for the little-$3$-cubes
operad on $\C^3 \simeq \R^6$, and the framing is supplied by the three choices
of holomorphic coordinate direction.

\noindent \emph{Caution:} the configuration-space topology $\mathrm{Conf}_2(\C^3)
\simeq S^5$ is simply connected, so the \emph{topological} braiding at the
$E_3$-level is trivial; the non-trivial content is in the $\dbar$-twisted
$L_\infty$-structure and in the Omega-background deformation
($h_1 + h_2 + h_3 = 0$), not in $\pi_1$.  This is the same situation as the
$E_2$-level on $\C^2$ (Remark~\ref{rem:en-from-dimension} in
quantum\_chiral\_algebras.tex): non-trivial braiding comes from the
Omega-background breaking, not from the topology of $\mathrm{Conf}_n$.

\subsection{Feynman computation: $1$-loop and $2$-loop coefficients}

Let $\mathbb{O}, \mathbb{O}'$ be two local operators at $z, w \in \C^3$,
both of classical dimension $1$ (say, components of $A$).  The OPE coefficient
in $A^{\hCS_6}_{\C^3}$ at order $\hbar^0$ is the classical BV bracket
\begin{equation}
  \{\mathbb{O}(z), \mathbb{O}'(w)\}_{\mathrm{cl}}
  \;=\; \frac{1}{(2\pi i)^3}\, \Omega_3 \wedge [\mathbb{O}, \mathbb{O}']_\frakg
  \cdot \delta^{(3)}(z - w) \cdot dz_1 \wedge dz_2 \wedge dz_3,
\end{equation}
a contact term in the Gelfand--Shilov distribution algebra on $\C^3$.

\noindent At $1$-loop, the first non-trivial Feynman integral involves the
bubble diagram with two $\dbar$-propagator insertions connecting two
$\mathrm{tr}([-, -] -)$ vertices.  The integrand is
\begin{equation}
  I^{\mathrm{bubble}}_1(z, w)
  \;=\; \int_{y \in \C^3} P(z, y) \wedge P(y, w)
  \wedge \mathrm{tr}(C_2^{\hCS_6}(\mathbb{O}, \mathbb{O}')),
\end{equation}
with $C_2^{\hCS_6}$ the $\frakg$-structure-constant tensor from the two
$[-,-]$-vertices.  The $P$-kernel squared integrates against $\Omega_3$
against $\dbar_y$-derivatives; a direct computation using the Fourier
representation $\widehat{P}(k) = i \bar{k}/\|k\|^2$ on the momentum side
gives
\begin{equation}
  I^{\mathrm{bubble}}_1(z, w) \;=\; \frac{1}{(2\pi i)^6}
  \cdot \frac{\log \|z - w\|^2}{\|z - w\|^6}
  \cdot \mathrm{tr}(T^a T^b T^a T^b),
  \label{eq:bubble-1loop}
\end{equation}
where $T^a$ are structure constants on $\frakg$.  The logarithm is the
one-loop beta-function signature of the Yangian Miura transformation;
the $\|z-w\|^{-6}$-scaling is Euler-homogeneous degree $-6$ on $\R^6$,
matching the expected CY-balanced short-distance asymptotic for a theory
on a complex $3$-fold (dimensional-analysis check: dimension of $A$ is
$[L]^0$; dimension of the product $A(z) A(w)$ is $[L]^0$; dimension of
the kinetic measure $\Omega_3 \wedge \dbar$ is $[L]^{-1}$; so the $6$d
propagator scales as $[L]^{-5}$ and the bubble as $[L]^{-10}$; the
logarithm is the dimensionless scaling anomaly).  For $\frakg = \fgl_1$,
$\mathrm{tr}(T^a T^b T^a T^b) = 1$, and we recover the Yangian Miura
$1$-loop coefficient $c_1 = 1$ predicted by Costello 2013 \S4
(Thm~\ref{thm:5d-boundary-yangian} in quantum\_chiral\_algebras.tex).

\noindent At $2$-loop, the leading non-trivial diagram is the theta
diagram ($3$ internal edges, $2$ trivalent vertices).  After BV
gauge-fixing and Euler-regularisation of the double $\dbar$-integrals,
the integrand evaluates to
\begin{equation}
  I^{\theta}_2(z, w) \;=\; \frac{1}{(2\pi i)^9}
  \cdot \frac{(\log \|z-w\|^2)^2}{\|z-w\|^6}
  \cdot c^{\hCS_6}_2,
  \label{eq:theta-2loop}
\end{equation}
with $c^{\hCS_6}_2 = (\mathrm{tr}\, \mathrm{ad}^2)^2 / 2 + \zeta(3)\,
(\mathrm{tr}\,\mathrm{ad})^3 / 6$ for semisimple $\frakg$; for
$\frakg = \fgl_1$ the second term vanishes and
$c^{\hCS_6}_2 = 1/2$.  The $\zeta(3)$ signals the standard CY$_3$
polylogarithmic content at $2$-loop; this matches the genus-$2$
Faber--Pandharipande BCOV / shadow-tower invariant in the chain-level
lane (cy\_to\_chiral.tex \S4645).

\subsection{$E_1$-chiral shadow on a reference curve}

Let $C \hookrightarrow \C^3$ be a holomorphic curve, say the $z_3$-axis
$C = \{z_1 = z_2 = 0\}$.  The factorisation pushforward (Costello--Gwilliam
\emph{FA} Vol 2 Prop 8.2.1) gives the $E_1$-chiral algebra on $C$:
\begin{equation}
  \Phi^{\mathrm{Sp}}_{C}\bigl(\Obs^q_{\hCS_6}\bigr) \;:=\;
  \int_{\C^2_{z_1, z_2}} \Obs^q_{\hCS_6} \;\Big|_{C}
  \;\in\; E_1\text{-ChirAlg}(C),
  \label{eq:e1-specialisation}
\end{equation}
where $\int_{\C^2_{z_1, z_2}}$ is transverse factorisation homology along
the $(z_1, z_2)$-plane at each point of $C$.  For $\frakg = \fgl_1$, this
specialisation is the chiral algebra $Y^+(\widehat\fgl_1)$ on $C$ (one
copy per point), which is exactly the CoHA-$\C^3$ route (a) and the
Costello--Li-$\C^3$ route (e) presentation of the CoHA.

\noindent For $C \subset K3 \times E$ a choice of elliptic fibre, the
corresponding specialisation (replacing $\C^2$ by $K3$) is
\begin{equation}
  \Phi^{\mathrm{Sp}}_{K3, E}\bigl(\Obs^q_{\hCS_6}(K3 \times E)\bigr)
  \;:=\; \int_{K3} \Obs^q_{\hCS_6} \;\Big|_E
  \;\in\; E_1\text{-ChirAlg}(E),
\end{equation}
and this object is the CY-A$_3$ chiral algebra $\Phi_3(D^b\Coh(K3 \times E))$
on $E$.  The five routes for $K3 \times E$ organise as five presentations
of the same factorisation-homology pushforward, not as five independent
functors.
\end{heal}

% ============================================================
\section{Heal 6 (consolidation): the CFG 2026 $E_3$-algebra of observables
  and the ``$\cW_{1+\infty}$'' identification}
\label{sec:ah6-cfg}
% ============================================================

\begin{theorem}[Costello--Francis--Gwilliam $E_3$-algebra of $\hCS_6$-observables]
\label{thm:cfg-e3}
Let $A^{\hCS_6}_{\C^3}$ be the cohomology of $\Obs^q_{\hCS_6}(\mathrm{Disk}_r(0))$
as $r \to 0$.  Then:
\begin{enumerate}[label=\textup{(\roman*)}]
  \item $A^{\hCS_6}_{\C^3}$ is an $E_3$-algebra in chain complexes over
    $\C[[\hbar]]$, with $E_3$-structure the three-fold nested-disk
    composition in $\C^3$.
  \item $A^{\hCS_6}_{\C^3}$ is filtered by polyvector degree, with associated
    graded $\mathrm{gr}\, A^{\hCS_6}_{\C^3} \simeq \PV^\ast(\C^3) \otimes
    U(\frakg)$ as $E_3$-algebras.
  \item For $\frakg = \fgl_1$, the $E_1$-specialisation to a curve
    $C \subset \C^3$ via~\eqref{eq:e1-specialisation} recovers the positive
    affine Yangian:
    \[
      \Phi^{\mathrm{Sp}}_C\bigl(A^{\hCS_6}_{\C^3}\bigr)
      \;\simeq\; Y^+(\widehat\fgl_1) \;=\; \mathrm{CoHA}(\C^3).
    \]
  \item The Drinfeld double completion is
    $\cD_\hbar(\Phi^{\mathrm{Sp}}_C(A^{\hCS_6}_{\C^3})) \simeq
    \cW_{1+\infty}|_{c=1}$, obtained by $\C^3 \sqcup (\C^3)^{\mathrm{op}}$
    gluing along $S^5$-link.
\end{enumerate}
\end{theorem}

\begin{proof}[Proof sketch with lane discipline]
\textbf{(i)} by Costello--Gwilliam \emph{FA} Vol 2 Thm~10.0.1 applied to
$\Obs^q_{\hCS_6}$: the pre-factorisation algebra of local observables of a
BV theory in $n$ complex dimensions is an $E_n$-algebra in cochain complexes
at each point.  For $\hCS_6$ on $\C^3$, $n = 3$.

\noindent \textbf{(ii)} filtration by polyvector degree: $\PV^{p,q}(\C^3)$
has $\dbar$-cohomology $H^{p,q}_{\mathrm{Dolb}}(\C^3; \wedge^p T_{\C^3}) =
\mathrm{Sym}^p \C^3 \otimes \mathrm{Sym}^q \bar\C^3$ for $\C^3$ affine; and
the $U(\frakg)$-tensor comes from the $\frakg$-coefficient of the BV fields.
The $E_3$-structure on the associated graded is the trivial braiding
composed with the $\PV$-Poisson bracket and the $U(\frakg)$-multiplication.

\noindent \textbf{(iii)} by Schiffmann--Vasserot 2017 \emph{Publ IHES} 118
Thm~A: the CoHA of the $\C^3$-one-loop quiver (equivalently, of $\PV^\ast(\C^3)
\otimes \fgl_1$) is the positive affine Yangian $Y^+(\widehat\fgl_1)$.  The
$E_1$-specialisation $\Phi^{\mathrm{Sp}}_C$ on the CFG $\hCS_6$ observables
produces the same algebra: Costello 2013 arXiv:1303.2632 \S4 shows the
$5$d-$\hCS$-boundary is $Y^+$ at $\fgl_1$, and the $5$d theory is the
$\C^3 \to \C^2 \times \R$ dimensional reduction of $\hCS_6$.

\noindent \textbf{(iv)} Drinfeld doubling: Arbesfeld--Schiffmann 2013
arXiv:1209.0429 Thm~2 constructs the doubling $\cD_\hbar(Y^+) = Y$; then
Prochazka--Rapcak 2019 gives $Y(\widehat\fgl_1) \simeq \cW_{1+\infty}|_{c=1}$.
\end{proof}

\begin{remark}[scope of the ``(a)-(e) coincidence'']
\label{rem:5routes-coincidence-scope}
Theorem~\ref{thm:cfg-e3}(iii) gives the clean statement of the ``five routes
coincide'' claim.  The coincidence is:
\begin{itemize}
  \item at $\hbar^0$ (classical associated graded) and $\hbar^1$ (first
    quantum correction), all five presentations agree by direct comparison.
  \item to all orders in $\hbar$, the coincidence is proved along the
    factorisation-homology pushforward spine; the chain-level equivalence
    is a theorem (Costello 2013 + Schiffmann--Vasserot 2017 + Costello--Li
    locality).
  \item the coincidence is of $E_1$-chiral algebras on $C$, \emph{not} of
    $E_3$-factorisation algebras on $\C^3$.  Different presentations exist
    for the $E_3$ object, and the choice of presentation matters for
    computing Wilson coefficients at higher loop order.
\end{itemize}
On $K3 \times E$: the five routes for $K3 \times E$ are five presentations
of the $E_1$-specialisation of $\Obs^q_{\hCS_6}(K3 \times E)$ along the
$K3$-fibre; they agree at the associated-graded level (Hodge
supertrace $\chi(\cO_{K3 \times E}) = 0$, spectral parameter
$u = 2\pi i \epsilon_1^E / \hbar$) and by the parallel Schiffmann--Vasserot
argument applied to $\Coh(K3)$ instead of $\Coh(\C^3)$.  Different $K3 \times
E$ \emph{different CY$_3$}-presentations (abelian-fibred vs. $K3$-fibred)
give different $E_1$-chiral shadows; this is AP-CY144 made concrete.
\end{remark}

% ============================================================
\section*{Lane discipline and cross-references}
% ============================================================

\begin{enumerate}[label=\textbf{L\arabic*.}]
  \item Five routes $= $ five distinct constructions, mutually
    $E_1$-quasi-isomorphic at the level of the factorisation-homology
    pushforward to a reference curve.  Not five applications of one
    functor.  (AP-CY144; AP-CY57; AP273.)
  \item Omega-background with independent $(\epsilon_1, \epsilon_2)$ is
    \emph{toric}-scoped; on $K3 \times E$ parameters collapse to the
    self-dual slice, and the surviving spectral parameter is $u = 2\pi i
    \epsilon_1^E / \hbar$.  (AP-CY20; AP17G.)
  \item Dunn additivity $E_2 = E_1 \otimes_{\mathsf{BV}} E_1$ is the Boardman--Vogt
    tensor; ``freezing'' means $S^1$-equivariant localisation, not a primitive
    operadic operation.  (AP16F; AP17E.)
  \item $\Obs^q_{\hCS_6}$ is a holomorphic pre-factorisation algebra on $\C^3$
    constructed from the BV superfield $\varphi = c + A + A^\dagger +
    c^\dagger$; the action $S_{\mathrm{cl}} = \int_{\C^3} \Omega_3 \wedge
    \mathrm{tr}(\tfrac{1}{2}\varphi \dbar \varphi + \tfrac{1}{6}\varphi
    [\varphi, \varphi])$ is a top-form after BV completion, resolving the
    ``$(3,2)$-form'' obstruction. (Costello \emph{Renormalization} Thm 8.4.1;
    Costello--Gwilliam \emph{FA} Vol 2 Thm 10.0.1; Costello--Li 2020
    arXiv:1505.06703; Costello--Francis--Gwilliam 2024.)
  \item $E_3$-algebra of observables on $\C^3$ specialises to $Y^+$ on $C$;
    doubling gives $\cW_{1+\infty}$; CoHA-$\C^3$ $=$ $Y^+$, not the full
    $\cW_{1+\infty}$.  (Cache 22T; Arbesfeld--Schiffmann 2013;
    Prochazka--Rapcak 2019.)
\end{enumerate}

\noindent The manuscript loci that should be updated in consequence:
\begin{itemize}
  \item \texttt{chapters/theory/quantum\_chiral\_algebras.tex} Route~(c) and
    Remark~\texttt{rem:6d-not-lagrangian}: link to the explicit BV
    construction~\eqref{eq:hcs6-Scl}--\eqref{eq:obs-6d}; state
    Thm~\texttt{thm:cfg-e3} in lieu of Conjecture~\texttt{conj:6d-boundary-toroidal}.
  \item \texttt{working\_notes.tex} L1665--1672: prefix the five-route list
    with the sentence ``five presentations of the $E_1$-specialisation of
    $\Obs^q_{\hCS_6}(\C^3)$ along the reference curve $C$'' and replace
    ``five independent'' by ``five distinct presentations.''
  \item \texttt{working\_notes.tex} L1694: replace ``freezes one of the two
    $E_1$-factors'' with ``$S^1_{\epsilon_2}$-equivariant localisation on
    the second factor, yielding $E_2^{S^1} \simeq E_1$'' and cite
    Nekrasov--Okounkov 2014.
  \item \texttt{chapters/theory/cy\_to\_chiral.tex} Thm~\texttt{thm:toric-chart-gluing}
    step~(iv): the Costello--Li comparison is the specialisation along the
    reference curve of $\Obs^q_{\hCS_6}$ (not of $5$d-$\hCS$ boundary), and
    the hocolim-of-envelopes identity is a consequence of the universal
    property of $\Phi^{\mathrm{Sp}}_C$.
  \item \texttt{chapters/frame/preface.tex} cascade table: add a column
    ``$d$-complex-dim holomorphic pre-factorisation algebra
    $\Obs^q_{\hCS_{2d}}$'' before the current ``$E_n$-chiral on $C$'' column
    and populate: $d = 1$ holomorphic on $E$; $d = 2$ on $K3$; $d = 3$ on
    $\C^3$ (and $K3 \times E$); $d = 5$ on $K3 \times K3 \times E$.
\end{itemize}

\end{document}
